% Physical unit 043; accepted Session 03 after fresh 400-dpi authority replay.
\GordanUnitStart{043}{printed page 38}
\GordanTopFolio{38}
\begingroup\small
Unde fit:
\[
\left\{\tg^2l_0''+\tg l_1\tg l_2\right\}\left\{\pi-\psi-e^2{\pi\over2}\cos l_0''\right\}^2=\left\{\tg l_1+\tg l_2\right\}^2
\]
\[
 \cos^2l_1\cos^2l_2\left\{{1\over\cos^2l_0''}-{\cos(l_1+l_2)\over\cos l_1\cos l_2}\right\}\left\{\pi-\psi-{e^2\over2}\pi\cos l_0''\right\}^2=\sin^2(l_1+l_2)
\]
\[
 \cos l_1\cos l_2\left\{{\cos l_1\cos l_2\over\cos^2l_0''}-\cos(l_1+l_2)\right\}\left\{\pi-\psi-{e^2\pi\over2}\cos l_0''\right\}^2=\sin^2(l_1+l_2)
\]
vel posito:
\[
\text{loco }\cos l_2:\ \cos l_1;\quad \text{loco }\cos(l_1+l_2):\ 1;\quad \text{loco }\sin(l_1+l_2):\ l_1+l_2,
\]
\[
 \cos^2l_1\left\{{\cos^2l_1\over\cos^2l_0''}-1\right\}\left\{\pi-\psi-e^2{\pi\over2}\cos l_0''\right\}^2=(l_1+l_2)^2
\]
\[
 \cos l_1\cot T\left\{\pi-\psi-{e^2\over2}\pi\cos l_0''\right\}=l_1+l_2.
\]
Qua in formula $T$ constat significare Azimut geodeticae lineae $l_0''$ in puncto $A$; ergo spectat in $l_0''$ formula:
\[
 \sin T\cos l_1=\cos l_0''.
\]
Posito hoc valore quantitatis $\cos l_0''$ in nostra aequatione, existit formula:
\[
 \cos l_1\cot T\{\pi-\psi\}-e^2{\pi\over2}\cos^2l_1\cos T=l_1+l_2
\]
e qua $T$ itemque $l_0''$ facile computari possunt.

Illa aequatio petitur ex his rationibus geometricis. --- Linea geodetica $l_2$, producta ultra punctum $B$, donec circulum \gq{$-l_0$} (e. g. in $C$) secet, cum circulo $-l_1$ et cum meridiano puncti $B$ (quas curvas in $D$ secantes facimus) efficit rectangulum triangulum, cujus latera tam brevia sunt, ut planus, possit triangulus intelligi et sumi:
\[
 {BD\over CD}=\tg BCD \qquad\text{cf. cap. I \S\ 9.}
\]
Atqui, ut facile intelligitur, est:
\[
 BD=a\int_{-l_1}^{l_2} dl\sqrt{1-e^2\cos^2l}
\]
\[
 CD=a\cos l_1\{2\mathcal A(l_0)-\psi\}
\]
et angulus $BCD$ complimentum ipsius $T$ efficit.
Quibus relatis in nostram formulam obtinemus:
\[
 \cot T\cos l_1\{2\mathcal A(l_0)-\psi\}=\int_{-l_1}^{l_2}dl\sqrt{1-e^2\cos^2l}
\]
vel neglecta excentricitate sub signo integrationis et in serie longitudinis periodi evoluta cap. III. tantummodo priorum duorum membrorum habita ratione,
\[
 \cot T\cos l_1\left\{\pi-e^2{\pi\over2}\cos l_0-\psi\right\}=l_1+l_2,
\]
id, quod erat demonstrandum.
\endgroup
